\chapter{Lesson 6: Asking Directions}

\section{Vocabulary}

\begin{center}
\begin{tabular}{llll}
\toprule
Characters & Pinyin & Part of speech & Meaning \\
\midrule
\zh{东} & dōng & noun & east \\
\zh{南} & nán & noun & south \\
\zh{西} & xī & noun & west \\
\zh{北} & běi & noun & north \\
\zh{上} / \zh{下} & shàng / xià & noun & above / below \\
\zh{前} / \zh{后} & qián / hòu & noun & front / back \\
\zh{左} / \zh{右} & zuǒ / yòu & noun & left / right \\
\zh{里} / \zh{外} & lǐ / wài & noun & inside / outside \\
\zh{旁边} & pángbiān & noun & beside \\
\zh{中间} & zhōngjiān & noun & middle, between \\
\zh{对面} & duìmiàn & noun & opposite \\
\zh{附近} & fùjìn & noun & nearby \\
\zh{离} & lí & preposition & (distance) from \\
\zh{远} & yuǎn & adjective & far \\
\zh{近} & jìn & adjective & near \\
\zh{往} & wǎng & preposition & towards \\
\zh{拐} & guǎi & verb & to turn \\
\zh{一直} & yìzhí & adverb & straight, continuously \\
\zh{路口} & lùkǒu & noun & intersection \\
\zh{红绿灯} & hónglǜdēng & noun & traffic light \\
\zh{过} & guò & verb & to cross, pass \\
\zh{米} & mǐ & measure word & metre \\
\zh{公里} & gōnglǐ & measure word & kilometre \\
\zh{地图} & dìtú & noun & map \\
\zh{迷路} & mí lù & verb-object & to get lost \\
\zh{走} & zǒu & verb & to walk, go \\
\zh{楼} & lóu & noun & building; floor \\
\zh{活动中心} & huódòng zhōngxīn & noun & activity centre \\
\bottomrule
\end{tabular}
\end{center}

\section{Direction and location words}

Location words are formed by adding \zh{边} \emph{bian} (or \zh{面} \emph{mian}) to
the basic directions: \zh{上边, 下边, 前边, 后边, 左边, 右边, 里边, 外边, 东边,
南边}. After a noun, the one-syllable forms \zh{上} and \zh{里} are common:
\zh{桌子上} ``on the table'', \zh{教室里} ``in the classroom''.

\begin{center}
\begin{tikzpicture}
  \draw[thick] (0,0) circle (2.2);
  \draw[->, very thick] (0,0) -- (0,2.0);
  \draw[->, very thick] (0,0) -- (0,-2.0);
  \draw[->, very thick] (0,0) -- (2.0,0);
  \draw[->, very thick] (0,0) -- (-2.0,0);
  \draw[->, thick, gray] (0,0) -- (45:1.9);
  \draw[->, thick, gray] (0,0) -- (135:1.9);
  \draw[->, thick, gray] (0,0) -- (-45:1.9);
  \draw[->, thick, gray] (0,0) -- (-135:1.9);
  \node[font=\large] at (0,2.55) {\zh{北}};
  \node[font=\large] at (0,-2.55) {\zh{南}};
  \node[font=\large] at (2.6,0) {\zh{东}};
  \node[font=\large] at (-2.6,0) {\zh{西}};
  \node at (45:2.75) {\zh{东北}};
  \node at (135:2.75) {\zh{西北}};
  \node at (-45:2.75) {\zh{东南}};
  \node at (-135:2.75) {\zh{西南}};
  \node[align=left, font=\small, anchor=west] at (3.6,0.6) {Chinese says \emph{east}/\emph{west} first:\\ \zh{东北} ``east-north'' = northeast,\\ \zh{西南} ``west-south'' = southwest.};
  \node[align=left, font=\small, anchor=west] at (3.6,-1.0) {The \zh{东北} region (Manchuria)\\ is literally ``the Northeast''.};
\end{tikzpicture}
\end{center}

\section{Grammar}

\subsection{Three ways to say where things are}

\begin{center}
\begin{tabular}{lll}
\toprule
Pattern & Focus & Example \\
\midrule
thing + \zh{在} + place & where a known thing is & \zh{图书馆在宿舍东边。} \\
place + \zh{有} + thing & what exists at a place & \zh{宿舍东边有一个图书馆。} \\
place + \zh{是} + thing & identifying what is there & \zh{宿舍东边是图书馆。} \\
\bottomrule
\end{tabular}
\end{center}

\begin{important}[\zh{在} vs \zh{有}]
\zh{在} takes a \emph{definite} subject (the library you both know about).
\zh{有} introduces an \emph{indefinite} object (``there is a library''). So
*\zh{一个图书馆在宿舍东边} is wrong; say \zh{宿舍东边有一个图书馆}.
\end{important}

\subsection{Distance with \zh{离}}

A \zh{离} B + \zh{远/近} or a distance.
\begin{itemize}
  \item \zh{我家离学校很近。} My home is close to school.
  \item \zh{机场离这儿有多远？}---\zh{大概二十公里。} How far is the airport from
        here?---About 20 km.
  \item \zh{离考试只有三天了。} There are only three days left until the exam.
\end{itemize}

\begin{important}[\zh{离} vs \zh{从}]
\zh{离} expresses static distance; \zh{从} expresses a starting point of motion:
\zh{从这儿走到那儿要十分钟}, but \zh{那儿离这儿不远}.
\end{important}

\subsection{Giving directions: \zh{往} + direction + verb}

\begin{itemize}
  \item \zh{一直往前走。} Go straight ahead.
  \item \zh{到第二个路口往右拐。} Turn right at the second intersection.
  \item \zh{过两个红绿灯，就到了。} Pass two traffic lights and you're there.
  \item \zh{往东走三百米。} Walk 300 metres east.
\end{itemize}

In southern China and Taiwan people tend to say \zh{左/右}; in Beijing, whose streets
form a grid, people often say \zh{东/西/南/北} even for walking directions.

\subsection{Asking the way}

\begin{itemize}
  \item \zh{请问，去火车站怎么走？} Excuse me, how do I get to the train station?
  \item \zh{请问，这附近有没有银行？} Is there a bank near here?
  \item \zh{坐几路车？} Which bus number?
  \item \zh{我迷路了。} I'm lost.
\end{itemize}

\section{Campus map}

The final oral task for this lesson used a map of an invented campus.

\begin{center}
\begin{tikzpicture}[b/.style={draw, thick, fill=#1, align=center, font=\small, minimum width=1.9cm, minimum height=1.2cm}]
  % roads
  \fill[gray!20] (-0.5,1.6) rectangle (11.5,2.4);
  \fill[gray!20] (4.6,-2.5) rectangle (5.4,5.5);
  \node[font=\scriptsize] at (0.3,2) {\zh{大学路}};
  \node[font=\scriptsize, rotate=90] at (5,-1.8) {\zh{中山路}};
  % light
  \fill[red] (5.65,2.55) circle (0.08); \fill[green!60!black] (5.85,2.55) circle (0.08);
  % buildings
  \node[b=blue!12] (lib) at (1.2,4) {\zh{图书馆}};
  \node[b=orange!15] (dorm) at (1.5,0) {\zh{宿舍}};
  \node[b=green!15] (gym) at (8,4) {\zh{体育馆}};
  \node[b=yellow!25] (cafe) at (3.4,4) {\zh{咖啡馆}};
  \node[b=purple!12] (cntr) at (8,0) {\zh{活动中心}};
  \node[b=red!12] (bank) at (10.3,0) {\zh{银行}};
  \node[draw, star, star points=5, fill=red!70, inner sep=2pt] (you) at (1.5,1.2) {};
  \node[font=\scriptsize, anchor=west] at (1.7,1.2) {\zh{你在这儿}};
  \draw[->, very thick, dashed, blue!70!black] (1.5,1.3) -- (1.5,2) -- (5,2) -- (5,1) -- (6.7,1) -- (6.8,0.65);
  % compass
  \draw[->, thick] (11,4) -- (11,5) node[above, font=\small] {\zh{北}};
\end{tikzpicture}
\end{center}

\begin{example}[Directions to the activity centre]
\zh{从宿舍出来，往北走到大学路，往东拐，一直走，到第一个路口（有红绿灯）往南拐，走五十米，活动中心在路的左边，银行旁边。}

Location descriptions from the map: \zh{咖啡馆在图书馆东边。图书馆对面是宿舍。活动中心和银行中间没有楼。体育馆离宿舍比较远。}
\end{example}

\begin{aside}
Check: if you walk south on Zhongshan Road, east is on your \emph{left}. The building
on the left is therefore on the east side of the road.
\end{aside}

\section*{Review}
\begin{itemize}
  \item Location words: direction + \zh{边/面}; noun + \zh{上/里}.
  \item \zh{东北} = northeast (east/west first).
  \item Definite thing \zh{在} place; place \zh{有} indefinite thing; place \zh{是} thing.
  \item A \zh{离} B \zh{远/近} for static distance; \zh{从} for starting points.
  \item \zh{往} + direction + \zh{走/拐}; \zh{一直往前走}.
\end{itemize}
